\section{Conclusion}\label{sec:conclusion}

The calibration procedure used in this paper provides accurate continuous-time
CARMA marginals for prices, temperature, solar production and wind production.
The ACF initialization gives interpretable persistence scales, while the QMLE
step improves the likelihood without losing the main autocorrelation structure.
Compared with discrete ARMA benchmarks, the CARMA specification keeps a
continuous-time state-space representation that is directly usable for joint
simulation.

For the coupled model, the relevant dependence is measured at the state-space
and driver level rather than on raw residual levels. Temperature and wind
transmit small but visible shocks to the price driver, with standard-deviation
ratios around \(4\%\), while the solar coupling is much weaker. This questions
the relevance of a linear coupled price--solar shock model on this dataset. More
generally, the present framework lets physical factors affect prices only
through shocks. Allowing temperature, wind or solar to also enter the price
drift could capture slower structural effects, and a multivariate CARMA
specification is a natural direction for such an extension.
